\section{Rational saturation and separated transport weights}
\label{sec:positive-new}

The two positive-domain companions below improve the wave-10 snapshot at
eleven and twelve unknowns. They use the complete compact source-44 tuple
\[
 (\beta,K_0,M_C,M_F),\qquad \beta=2^{590069}-1,
\]
with the already shifted high mask. The tuple must not be mixed with the
source-34 tuple used by some integer-domain constructions. Its inherited
native compiler and syntactic-refutation calibration remain dependencies.
The new work is an exact existence reduction on the normalized positive
Pell source, rather than a new claim about those compiler theorems.

The changes compose three mechanisms which were separate in wave 10:
source 69's rational Pell saturation, source 70's three-root transport,
and source 74's exact treatment of the possible common factor three in
the adaptive main modulus. A fourth observation separates a polynomial
which bounds signed roots from the multiplier which provides a positive
inverse transport coordinate. They need not be the same polynomial.

\begin{center}
\begin{tabular}{lrrrr}
\hline
Equation & Positive unknowns & Exact degree & Gates & System gates\\
\hline
Source 70 & 11 & 325 & 113 & 110\\
$P_{11}^{\mathrm{new}}$ & 11 & 269 & 111 & 108\\
Source 69 & 12 & 82 & 102 & 96\\
$P_{12}^{\mathrm{new}}$ & 12 & 78 & 102 & 96\\
\hline
\end{tabular}
\end{center}
The eleven-variable result reduces degree by $56/325$, approximately
$17.23\%$, and removes two arithmetic gates. The twelve-variable result
reduces degree by $4/82$, approximately $4.88\%$, at unchanged cost.
Every live binary addition, subtraction, and multiplication counts as one
gate, including multiplication by a fixed integer. Polynomial register
names do not add quantified variables.

\subsection{Rational Pell saturation}

We reproduce the elementary source-69 theorem because it removes an
entire divisibility restriction from the eleven-variable construction.

\begin{lemma}[Rational-coordinate saturation]
\label{pp:saturation}
If $Q>1$ and $V$ are integers and $y$ is rational, then
\[
 QV^2-(Q-1)y^2=1\quad\Longrightarrow\quad y\in\mathbb Z.
\]
\end{lemma}
\begin{proof}
The equation excludes $V=0$ and $y=0$. Change signs to assume both are
positive. Write $y=a/m$ in lowest terms, $m>0$. Clearing denominators
shows $m^2\mid Q-1$, so $(Q-1)y$ is integral. The map
\[
 V'=(2Q-1)V-2(Q-1)y,\qquad y'=(2Q-1)y-2QV
\]
preserves the equation and makes $V'$ integral. When $V>1$, the identities
\[
 y^2-V^2=\frac{V^2-1}{Q-1}>0,
 \qquad (QV)^2-((Q-1)y)^2=QV^2+Q-1>0
\]
give $-V<V'<V$. The norm excludes $V'=0$. Moreover, $Q\equiv1\pmod m$
implies that the numerator of $y'$ over denominator $m$ is congruent to
$a$ modulo $m$. Its reduced denominator is still exactly $m$.
Replace the new coordinates by their absolute values and repeat. The
positive integral first coordinate strictly decreases until it is one,
where the equation forces $y^2=1$ and hence $m=1$. Since the reduced
denominator was preserved, the initial $y$ was integral.
\end{proof}

This is an unbounded descent argument. It requires $Q$ and $V$ integral;
it cannot be replaced by a claim that every rational Pell coordinate is
integral. For example, $2\cdot1^2-4\cdot(1/2)^2=1$ shows why the adjacent
coefficients $Q$ and $Q-1$ matter.

\subsection{The weighted source interface and its adaptive modulus}

The inherited normalized source supplies the following registers:
\begin{align*}
 q&=\beta J+1,& C&=q-F-\alpha,&E&=XY,\\
 k&=\eta+\zeta,&c&=kY+\eta,&a&=E+Y,\\
 H&=4a+3,&\Delta&=a^2+H=(a+2)^2-1,&\ell&=EkY,\\
 R&=((q-1)(q-F)+\alpha)(q^2-1)+(M_C+qM_F)J.
\end{align*}
Here the supplied witnesses are strictly positive. The two scale
divisibilities and two deleted quotient conditions are restored by
carriers. Given any positive polynomial value $h$, define
\begin{align}
 W&=1+(q-1)h,&M&=q^3(q-1),\nonumber\\
 B_t&=(qK_0+X)C-qF,\nonumber\\
 Z&=M(\ell+k+r)+(q-1)Y-q^3WB_t,\nonumber\\
 A_0&=Z^2-M^2\{\ell(\ell+k)+1\},                 \label{pp:first}\\
 S&=qY,&U&=qS=q^2Y,&v&=k-R-1,\nonumber\\
 D_*&=U(X+ac)+H(U\gamma+SX-v),\nonumber\\
 A_1&=D_*^2-U^2(1+\Delta c^2).                  \label{pp:main}
\end{align}
The other source equations are the normalized strong and auxiliary norms
\begin{align}
 f^2-\Delta i^2c^4&=1,\nonumber\\
 QV^2-(Q-1)y^2&=1,\qquad
 Q=\Delta(f^2-1),\quad V=f^2(cT-R)-c.            \label{pp:strongaux}
\end{align}
The signed auxiliary coordinate is harmless; after recovery its positive
gap is $|y|-V>0$. The strong root is replaced by its absolute value when
needed.

For clarity, the exact index-recovery argument using the smaller modulus
$U=q^2Y$ is recalled here. From $A_0=0$, the rational root $Z/M$ has an
integer square and is integral. Reduction modulo $q^3$ and $q-1$ gives
\[
 q^3\mid Y,\quad q-1\mid B_t,\quad
 \tau=\ell+k+r+s-Wt,\quad
 s=Y/q^3>0,\quad t=B_t/(q-1).
\]
Since $|\tau|<\ell+k$, we have $Wt>r+s>0$. Thus $B_t>0$ and $C>0$.
The authentic mask bounds now give $q\ge16$ and $3q<R<q^4-2$.
The main numerator is positive, since
$D_*>Uac-Hk>k(UaY-H)>0$. Rational-square integrality applied to
$A_1=0$ gives a positive integral main Pell root.

Write $A=a+2$ and $P_0=2XY^2+1$. The inherited Pell classification gives
indices $n_0,p$ with
\[
 |\tau|=\chi_{P_0}(n_0),\quad k=2\psi_{P_0}(n_0),\qquad
 D_*/U=\chi_A(p),\quad c=\psi_A(p).
\]
It gives $n_0<p$, $n_0\ge2$, $c>8XY^3\ge8q^9>2R$, and $c>2p$.
The normalized strong norm makes $f=\chi_A(m)$ with $c\mid m$ and
$m>2p$. For example, strong divisibility first gives $p\mid m$; writing
$m=pv_0$ and reducing $\psi_{\chi_A(p)}(v_0)$ modulo $c$ shows
$c\mid v_0$ from $c^2\mid ic^2$. The auxiliary congruences
$V\equiv-c\pmod f$ and $V\equiv-R\pmod c$, together with its norm,
then invoke the inherited signed auxiliary-index lemma to yield $p=R$.
None of these steps assumes the deleted scale $q\mid X$.

In both new constructions $J\equiv1\pmod3$; since $\beta\equiv1\pmod3$,
$q\equiv2\pmod3$. Therefore $\gcd(q,H)=1$ and
\[
 g:=\gcd(H,U)=\gcd(3,Y)\in\{1,3\}.
\]
Crucially, $q^3\mid Y$ and $\gcd(q,g)=1$ imply $q^3\mid Y/g$. Hence
\[
 \mu=\frac{qY}{g}\ge q^4>R.
\]
The congruence $U\mid H(SX-v)$ gives $\mu\mid v$. The first recurrence
gives $k\equiv2n_0\pmod{4XY^2}$, and $\mu\mid4XY^2$. Thus
$\mu\mid2n_0-R-1$. The strict index bounds force $2n_0=R+1$.
Consequently $v=k-2n_0$ is a positive multiple of $4XY^2$, and so of $U$.
Only now does the main congruence give $q\mid HX$, and then $q\mid X$.
The recovered source coordinates are
\[
 w=X/q>0,\quad h_{\rm index}=v/(XY)>0,\quad
 \gamma_{\rm old}=\gamma+w-v/U>0.
\]
The last inequality follows from
$H\gamma_{\rm old}=D_*/U-ac-X>c-X>0$.
The first carrier restores the positive transport quotient. This is
the same normalized compiler source as in the wave-10 packages.

Conversely its canonical witnesses have $q=B^N$, $B=\beta+1$, with
odd $N>1$, $R\equiv3\pmod4$, $X=2^R$, $q^3\mid Y$, and
\[
 t_0=B_t/(q-1)=q(z_{\rm source}-1)\ge q,\quad
 s=Y/q^3<k,\quad \tau>\ell.
\]
For $h\ge k$ the choice
\[
 r=Wt_0+\tau-\ell-k-s>qW-2k>0
\]
restores the first carrier. The main coordinate
$\gamma=\gamma_{\rm old}-w+v/U$ is positive, because
$v/U\ge4XY/q^2>w$. These are the precise completeness facts used below.

\subsection{The eleven-variable polynomial}

Supply exactly
\[
 j,F,\alpha,i,T,X,Y,n,\eta,\zeta,\gamma\in\mathbb Z_{>0},
 \qquad J=3j+1.
\]
Use the common registers above and put
\begin{align}
 I&=ic^2,&D_f&=1+\Delta I^2,&Q&=\Delta(D_f-1),\nonumber\\
 \delta&=Q-1,&V&=D_f(cT-R)-c,&D_y^{\rm num}&=QV^2-1,\nonumber\\
 \widehat R&=q^3(q+F+\alpha),&\overline V&=D_f(cT+\widehat R),
       &L&=3\overline V,\nonumber\\
 h&=D_fcT,&W&=1+(q-1)h,\nonumber\\
 \overline Z&=M(\ell+k+n+L)+(q-1)Y-q^3WB_t.       \label{pp:11regs}
\end{align}
Use $A_1$ from \eqref{pp:main} and define the octic norm schedule
\begin{align}
 \mathcal A&=M^2\{\ell(\ell+k)+1\},&
 \mathcal B&=M^2D_f,&\mathcal C&=M^2D_y^{\rm num},\nonumber\\
 T_n&=\delta^2\mathcal A\mathcal B,&R_n&=\delta\mathcal C,\nonumber\\
 J_n&=\delta(\overline Z^2-\mathcal A-\mathcal B)-\mathcal C,\nonumber\\
 N_*&=J_n^2-4\{T_n+R_n(\mathcal A+\mathcal B)\},\nonumber\\
 N_8&=N_*^2-64T_nR_n\overline Z^2.              \label{pp:11norm}
\end{align}
The requested ordinary polynomial is
\begin{equation}
 \boxed{P_{11}^{\rm new}=N_8+A_1+N_8A_1.}       \label{pp:11poly}
\end{equation}
There are no divisions in this definition. Its literal evaluator is
\texttt{p11\_positive\_degree269.json}.

\begin{theorem}
The polynomial \eqref{pp:11poly} has a strictly positive zero exactly when
the inherited normalized compiler source has a positive solution.
\end{theorem}
\begin{proof}
The equation is $(1+N_8)(1+A_1)=1$. As $N_8$ is a square modulo four,
$N_8\ne-2$. The two integer factors therefore must both be one. Thus the
single equation is equivalent to $N_8=A_1=0$. No parity restriction on
$A_1$ is needed for this two-factor step.

For the root extraction write
$D_\tau=\ell(\ell+k)+1$ and $D_y=D_y^{\rm num}/\delta$ in the proof only.
The standard eight-conjugate identity gives
\[
 N_8=\delta^4\prod_{\epsilon_1,\epsilon_2,\epsilon_3=\pm1}
 \left(\overline Z-M\epsilon_1\sqrt{D_\tau}
                   -M\epsilon_2\sqrt{D_f}-M\epsilon_3\sqrt{D_y}\right).
\]
Its radicands are positive before any equation is imposed. Indeed
$1<D_\tau<D_f<Q$; also, writing $b=cT-R\in\mathbb Z$, either $b=0$
and $V=-c$, or $|V|\ge D_f-c>1$. Therefore
\[
 D_y=V^2+\frac{V^2-1}{Q-1}>V^2>1.
\]
It cannot equal $D_\tau$ or $D_f$: equality to an integer $D$ with
$1<D<Q$ would give $Q(D-V^2)=D-1$, an impossibility. Thus the three
radicands are distinct.

Their signed sum cannot be zero. With positive roots $t,f,y$, one has
$t<(a+1)c$ and $f>(a+1)c^2$. If $b=0$, then $y<c+1$ and $f>t+y$.
If $b\ne0$, then $y>|V|\ge f^2-c>2f>f+t$.
Consequently a zero of the norm supplies signed roots with nonzero
rational sum
\[
 \overline Z/M=\tau_*+f_*+y_*.
\]
For completeness, a nonzero rational sum of three signed roots of
distinct positive rational numbers forces all three roots rational.
To see this, let $d_i=x_i^2$ and $t=x_1+x_2+x_3\ne0$. Then
\[
 u=\frac{t^2-d_1-d_2-d_3}{2},\qquad
 v=\frac{u^2-d_1d_2-d_1d_3-d_2d_3}{2t}
\]
are rational and equal respectively to the pairwise sum and product of
the roots. The identities
\[
 d_1t+v=x_1(x_1+x_2)(x_1+x_3),\qquad
 x_1=\frac{d_1(d_1+u)}{d_1t+v}
\]
recover $x_1$ rationally. The denominator is nonzero by distinctness of
the positive $d_i$, and the cyclic formulas recover the other roots.

Since $D_\tau,D_f$ are integers, $\tau_*,f_*$ are integers. The third
root satisfies $QV^2-(Q-1)y_*^2=1$, with integral $Q>1$ and $V$.
Lemma~\ref{pp:saturation} makes $y_*$ integral as well. This replaces
source 70's coprimality argument and permits $I=ic^2$ throughout.

The authentic mask inequalities $0<M_C<\beta$, $0<M_F<2\beta$ imply
$|R|<\widehat R$. Thus $|V|<\overline V$, $f<\overline V$, and
$y<2|V|<2\overline V$, so $L>f+y$. Set
\[
 r=n+L-f_*-y_*>0,\quad f_{\rm old}=|f_*|,\quad
 y_{\rm old}=|y_*|,\quad u=y_{\rm old}-V>0.
\]
The first carrier becomes $Z=\overline Z-M(f_*+y_*)=M\tau_*$.
The strong and auxiliary source norms hold by construction. All the
weighted source witnesses are now positive, and $h=D_fcT>0$.
The preceding adaptive argument restores the normalized source.

Conversely choose canonical normalized source witnesses and its positive
auxiliary reconstruction, with $f^2=D_f$, $y=V+u>0$. The latter can be
constructed from the odd index $R\equiv3\pmod4$ using
$y=\psi_{\Delta I}(R)$ and $V=\chi_{\Delta I}(R)/(\Delta I)$.
The squared-modulus congruences make
$T=(V+c+Rf^2)/(cf^2)$ a positive integer. No extra factor $q(q-1)$ in
$I$ is required. Here $F+\alpha<q$ and $c>8q^9$, so
\[
 \widehat R<2q^4<c\le cT,\qquad L<6D_fcT=6h.
\]
Also $h>k$. For the canonical first carrier,
\[
 r>qW-2k>q+\{q(q-1)-2\}h>6h>L.
\]
Thus $n=r-L+f+y>0$, and substitution gives
$\overline Z=M(\tau+f+y)$, so $N_8=0$. The adaptive main coordinate is
positive by the earlier completeness argument. Finally the canonical
repunit has $J\equiv1\pmod3$ and $J>1$, so $j=(J-1)/3>0$.
\end{proof}

\subsection{The twelve-variable polynomial}

Supply exactly
\[
 j,F,\alpha,t,T,X,Y,r,\eta,\zeta,n,\gamma\in\mathbb Z_{>0},
 \qquad J=6j+1.
\]
Use \eqref{pp:first} with $h=k$ and \eqref{pp:main}. Put
\begin{align}
 f&=t+1,&\rho&=f^2-1,&d&=\Delta c^4,\nonumber\\
 Q&=\Delta\rho,&V&=f^2(cT-R)-c,\nonumber\\
 a_*&=(V+n)^2-Q\{(V+n)^2-V^2\}-1,\nonumber\\
 b_*&=2(1-Q)(V+n),&A_*&=da_*+(1-Q)\rho,\nonumber\\
 N&=A_*^2-db_*^2\rho.                          \label{pp:12norm}
\end{align}
Define
\begin{equation}
 \boxed{P_{12}^{\rm new}=(1+A_0)(1+A_1)(1+N)-1.}\label{pp:12poly}
\end{equation}
This is precisely source 69's saturated merger with the smaller adaptive
main modulus. The literal evaluator is
\texttt{p12\_positive\_degree78.json}.

\begin{theorem}
The polynomial \eqref{pp:12poly} has a strictly positive zero exactly when
the inherited normalized compiler source has a positive solution.
\end{theorem}
\begin{proof}
Here $q$ is even, since $\beta$ and $J$ are odd. Therefore $A_0$ and
$A_1$ are squares modulo four; so is $N$, because $b_*$ is even.
The product equals one exactly when all three residuals vanish.

The first and main carriers recover $C>0$ and $c>2R$ before the strong
norm is used. Since $f\ge2$ and $T\ge1$,
$V=f^2(cT-R)-c>c>1$. Thus $b_*\ne0$ at a candidate zero.
Define rational candidates $i=A_* /(db_*)$ and $u=n-i$. The norm gives
$di^2=\rho$. If
\[
 H_0(v)=(1-Q)v^2+2(1-Q)Vv+V^2-1,
\]
then direct expansion gives $H_0(n-i)=0$. Hence $y=V+u$ satisfies
$QV^2-(Q-1)y^2=1$. Lemma~\ref{pp:saturation} gives $y,u,i\in\mathbb Z$.
Replace $i$ by $|i|>0$ and the gap by $|y|-V>0$; the positivity of
the latter follows from $y^2-V^2=(V^2-1)/(Q-1)>0$.
These are the normalized strong and auxiliary source witnesses, and the
adaptive carrier restores the remaining source quotients.

Conversely, canonical normalized witnesses supply positive $f,i,u$;
take $t=f-1>0$ and $n=i+u>0$. The identity
\[
 \left.N\right|_{\rho=di^2}
   =d^2 H_0(n-i)H_0(n+i)
\]
gives $N=0$. The canonical first and adaptive main carriers vanish by
the preceding reconstruction. The repunit has $J\equiv1\pmod6$ and
$J>1$, so $j=(J-1)/6>0$. This proves completeness.
\end{proof}

\subsection{Exact degrees, coefficient certificates, and audit scope}

Write $q_0=3\beta j$ for the eleven-variable polynomial and $q_0=6\beta j$
for the twelve-variable one. In either case put
\[
 k_0=\eta+\zeta,\quad R_4=q_0^3(q_0-F),\quad
 B_{t,2}=(K_0q_0+X)(q_0-F-\alpha)-q_0F.
\]
The cancellation in the main Pell norm gives
\[
 [A_1]_{13}=8q_0^2k_0X^2Y^4(q_0^2Y\gamma+R_4).
\]
For eleven variables, the exact degrees of the principal registers are
\[
 \deg I=5,\quad \deg D_f=14,\quad \deg Q=18,\quad
 \deg h=17,\quad \deg L=18,\quad \deg\overline Z=23.
\]
If $D_{f,14}=i^2k_0^4X^2Y^6$ and $Q_{18}=i^2k_0^4X^4Y^8$, then
\[
 [\overline Z]_{23}=-q_0^4D_{f,14}k_0YT B_{t,2},\qquad
 [N_8]_{256}=Q_{18}^4[\overline Z]_{23}^{\,8}.
\]
Multiplication by the displayed nonzero degree-13 main leader proves
the upper bound 269 is attained. More specifically, the coefficient of
\[
 j^{36}i^{24}\eta^{57}X^{42}Y^{93}T^8\alpha^8\gamma
\]
is exactly $8(3\beta)^{36}>0$. The value is independent of $K_0,M_C,M_F$,
so it remains nonzero after their actual fixed specialization.

For twelve variables the residual degrees are $15,13,50$. Their leading
forms are the main leader above together with
\[
 [A_0]_{15}=-2q_0^8k_0^2XY^2B_{t,2},\qquad
 [N]_{50}=4k_0^8X^8Y^{16}t^8R_4^2n^2.
\]
The coefficient of
$j^{18}F^2\eta^{11}X^{12}Y^{23}t^8n^2\gamma\alpha$ in
$P_{12}^{\rm new}$ is $64(6\beta)^{18}>0$. Thus its degree is exactly 78.

The included standard-library audits independently expand the literal
polynomial formulas, prove the universal norm and descent identities,
check all live inputs and gates, and compute these exact leading forms.
The twelve-variable audit checks 13,128 terms in the first carrier and
4,417 in the main carrier; it handles the merger as an exact 32-term
polynomial cut. These finite computations do not decide consistency or
replace the unbounded proofs. Composing the two existence equivalences
with the fixed source calibration, and negating, yields the corresponding
ZFC-provable no-positive-zero equivalences with $\operatorname{Con}(\mathrm{ZFC})$.
